Theorem~\ref{fubini} and Proposition~\ref{prop:geometric:fubini} allow
fibre representatives in $\RV[*]$, that is, finite disjoint unions of
objects of several degrees. Their proofs above fix one label degree for
each representative, with nonzero labels. In this section we treat the
general case. The bounded case needs no new idea; the geometric case
needs a padding of the vanishing labels, since the geometric value of a
summand of degree $k$ is normalized by $L^{-k}$.

\begin{defn}\label{defn:mixed:presentation}
Let $\mathbb T$ be a parameter base and $k\geq0$. A presentation in
degree at most $k$ is a pair $\bm U=(U,p)$, where
$U\subseteq\RV_\infty^a$ is $\mathbb T$-definable and
\[
 p:U\longrightarrow\RV_\infty^k
\]
is $\mathbb T$-definable and finite-to-one. For $I\subseteq\{1,\ldots,k\}$,
the exact stratum
\[
 U_I=\{u\in U:\ p_j(u)\neq\infty\ \Longleftrightarrow\ j\in I\},
\]
with the label $(p_j)_{j\in I}$, is an object $\bm U_I\in\RV[|I|]$.
The lift of $\bm U$ is
\[
 \mathbb L\bm U=\{(x,u)\in K^k\times U:\ \rv(x)=p(u)\}.
\]
Since $\rv(x_j)=\infty$ precisely when $x_j=0$, the lift is the disjoint
union over $I$ of copies of $\mathbb L\bm U_I$, with the coordinates
outside $I$ equal to zero. Conversely, a finite tagged disjoint union of
objects of degrees $k_1,\ldots,k_m$ has a presentation in degree at most
$\max_j k_j$, obtained by adding zero labels; its exact strata are the
given summands, and padding by zero coordinates identifies the lifts. This
reindexing is how statements for representatives in $\RV[*]$ follow from
statements for presentations; the formalization works with presentations
throughout and does not state the reindexing separately. Put
\begin{align}
 \mathbb E_{\mathrm b}^{\mathbb T}([\bm U])
   &=\sum_I\mathbb E_{\mathrm b}^{\mathbb T}([\bm U_I])\quad\text{in }R(\mathbb T),
 \label{mixed:bounded:value}\\
 \mathbb E_{\mathrm g}^{\mathbb T}([\bm U])
   &=\sum_I L^{-|I|}\,\widetilde{\mathbb E}_{\mathrm g}^{\mathbb T}([\bm U_I])
   \quad\text{in }R_g(\mathbb T).
 \label{mixed:geometric:value}
\end{align}
Thus each exact stratum is normalized in its own degree.
\end{defn}

The bounded value \eqref{mixed:bounded:value} is the bounded value of the
underlying RV set $U$: as in the proof of Theorem~\ref{fubini}, the bounded
retraction does not depend on the finite-to-one label. If all labels of
$\bm U$ are nonzero, then $\bm U=\bm U_{\{1,\ldots,k\}}$ is an object of
$\RV[k]$, and both values are the values of Section~\ref{sec:setup}.

\begin{lem}\label{lem:mixed:lift}
Let $\bm U$ be an $\mathbb S$-definable presentation in degree at most $k$.
Then $\mathbb L\bm U\in\VF_*$ and
\[
 \mathbb E_{\mathrm b}^{\mathbb S}\!\left(\int[\mathbb L\bm U]\right)
   =\mathbb E_{\mathrm b}^{\mathbb S}([\bm U]),
 \qquad
 \mathbb E_{\mathrm g}^{\mathbb S}\!\left(\int[\mathbb L\bm U]\right)
   =\mathbb E_{\mathrm g}^{\mathbb S}([\bm U]).
\]
\end{lem}

\begin{proof}
The set $\mathbb L\bm U$ is an RV pullback: membership of $(x,u)$
depends on $x$ only through $\rv(x)$. Its RV fibres are finite since $p$
is finite-to-one. Partition it according to the set $I$ of nonzero VF
coordinates. The piece indexed by $I$ is definably bijective, by deleting
the zero coordinates, to $\mathbb L\bm U_I$, and its key
$(\rv(x),u)\mapsto u$ identifies its RV data with $U_I$. Hence its integral
is the class of $\bm U_I$ modulo the special relation. Applying the two
retractions and adding over $I$ gives \eqref{mixed:bounded:value} and
\eqref{mixed:geometric:value}.
\end{proof}

Let $\bm{\mathbb A}_{\RES}\in\RES[1]$ be the disjoint union of $\G_m$ and the
degree-one point, as in the proof of Proposition~\ref{prop:geometric:fubini}.
Its unnormalized geometric value is $g+1=L$, and its lift is definably
bijective to $\mathcal O$.

\begin{defn}\label{defn:thickening}
Let $\bm U=(U,p)$ be a presentation in degree at most $k$, and let
$1\leq i\leq k$. Put $U^{\neq}=\{u\in U:p_i(u)\neq\infty\}$ and
$U^{=}=\{u\in U:p_i(u)=\infty\}$. The thickening of $\bm U$ at $i$ is
the tagged disjoint union
\[
 \theta_i\bm U=U^{\neq}\ \sqcup\ \bigl(U^{=}\times\bm{\mathbb A}_{\RES}\bigr),
\]
with the label $p$ on the first summand, and with the label of $(u,w)$ on
the second summand obtained from $p(u)$ by replacing its $i$th coordinate
$\infty$ by the label of $w$. Put
$\theta\bm U=\theta_k\cdots\theta_1\bm U$.
\end{defn}

\begin{lem}\label{lem:thickening}
Let $\bm U$ be a presentation in degree at most $k$ over a parameter base
$\mathbb T$, and let $1\leq i\leq k$.
\begin{enumerate}
\item The thickening $\theta_i\bm U$ is a presentation in degree at most $k$
over $\mathbb T$. It is obtained from the defining formulas of $\bm U$ by a
fixed first-order operation, uniformly in parameters, and a bound for its
label fibres depends only on a bound for those of $\bm U$ and a fixed
bound for those of $\bm{\mathbb A}_{\RES}$. All labels of
$\theta\bm U$ are nonzero, so $\theta\bm U\in\RV[k]$.
\item $\mathbb E_{\mathrm g}^{\mathbb T}([\theta_i\bm U])
       =\mathbb E_{\mathrm g}^{\mathbb T}([\bm U])$; hence
$\mathbb E_{\mathrm g}^{\mathbb T}([\theta\bm U])
       =\mathbb E_{\mathrm g}^{\mathbb T}([\bm U])$.
\item Let $Z\subseteq K^d\times\RV_\infty^q$ be $\mathbb S$-definable and let
$(\bm U_c)_{c\in Z}$ be an $\mathbb S$-definable family of presentations in
degree at most $k$ whose total lift
\[
 \mathbb L_Z\bm U=\{(c,x,u):c\in Z,\ (x,u)\in\mathbb L\bm U_c\}
\]
lies in $\VF_*$. Then $\mathbb L_Z\theta_i\bm U\in\VF_*$ and
\[
 \mathbb E_{\mathrm g}^{\mathbb S}\!\left(\int[\mathbb L_Z\theta_i\bm U]\right)
   =\mathbb E_{\mathrm g}^{\mathbb S}\!\left(\int[\mathbb L_Z\bm U]\right).
\]
\end{enumerate}
\end{lem}

\begin{proof}
For (1), the defining condition of $\theta_i\bm U$ is a disjunction of
two conditions: either $(u,p(u))$ is a point of $\bm U$ with $p_i(u)\neq\infty$
and the padding coordinates are zero, or the point obtained by setting
the $i$th label to $\infty$ lies in $\bm U$ and the padding coordinates
together with the new $i$th label lie on the label graph of
$\bm{\mathbb A}_{\RES}$. If $N$ bounds the label fibres of $\bm U$ and $N'$
those of $\bm{\mathbb A}_{\RES}$, then $N+NN'$ bounds those of
$\theta_i\bm U$. The label of the second summand at $i$ is nonzero, and
$\theta_i$ does not change the other labels. Hence $\theta\bm U$ has nonzero
labels.

For (2), the exact strata of $\theta_i\bm U$ are
\[
 (\theta_i\bm U)_J=\bm U_J\ \sqcup\ \bigl(\bm U_{J\setminus\{i\}}\times
 \bm{\mathbb A}_{\RES}\bigr)\quad(i\in J),
 \qquad (\theta_i\bm U)_J=\varnothing\quad(i\notin J).
\]
Since $\widetilde{\mathbb E}_{\mathrm g}$ is multiplicative and
$\widetilde{\mathbb E}_{\mathrm g}([\bm{\mathbb A}_{\RES}])=L$,
\[
 \mathbb E_{\mathrm g}([\theta_i\bm U])
 =\sum_{J\ni i}L^{-|J|}\widetilde{\mathbb E}_{\mathrm g}([\bm U_J])
  +\sum_{I\not\ni i}L^{-|I|-1}L\,\widetilde{\mathbb E}_{\mathrm g}([\bm U_I])
 =\mathbb E_{\mathrm g}([\bm U]).
\]

For (3), partition $\mathbb L_Z\bm U$ according to whether the $i$th lifted
coordinate $x_i$ vanishes. On the piece $x_i\neq0$, forgetting the zero
padding coordinates is a bijection onto the first summand of
$\mathbb L_Z\theta_i\bm U$. On the piece $x_i=0$, take its product with the
lift of $\bm{\mathbb A}_{\RES}$; reinserting the lifted padding coordinate
in position $i$ is a bijection onto the second summand. Both maps are
definable over $\mathbb S$, so $\mathbb L_Z\theta_i\bm U\in\VF_*$. Since
$\mathbb E_{\mathrm g}^{\mathbb S}\circ\int$ is a ring homomorphism and
$\mathbb E_{\mathrm g}^{\mathbb S}(\int[\mathbb L\bm{\mathbb A}_{\RES}])
=\mathbb E_{\mathrm g}^{\mathbb S}([\bm{\mathbb A}_{\RES}])
=L^{-1}\widetilde{\mathbb E}_{\mathrm g}([\bm{\mathbb A}_{\RES}])=L^{-1}L=1$, the two pieces
retain their geometric values.
\end{proof}

\begin{thm}\label{thm:mixed}
Let $\mathbb S$, $A,A'\in\VF_*$, $f:A\to B$, $f':A'\to B$, $A_b$, $A'_b$ and
$\mathbb S_b$ be as in Theorem~\ref{fubini}. Suppose that for every
$b\in B$ there are $k_b,l_b\geq0$, $\mathbb S_b$-definable presentations
$\bm U_b$ in degree at most $k_b$ and $\bm U'_b$ in degree at most $l_b$, and
$\mathbb S_b$-definable bijections
\[
 A_b\xrightarrow{\ \sim\ }\mathbb L\bm U_b,
 \qquad
 A'_b\xrightarrow{\ \sim\ }\mathbb L\bm U'_b .
\]
\begin{enumerate}
\item If
$\mathbb E_{\mathrm b}^{\mathbb S_b}([\bm U_b])
 =\mathbb E_{\mathrm b}^{\mathbb S_b}([\bm U'_b])$
in $R(\mathbb S_b)$ for every $b\in B$, then
$\mathbb E_{\mathrm b}^{\mathbb S}(\int[A])
 =\mathbb E_{\mathrm b}^{\mathbb S}(\int[A'])$ in $R(\mathbb S)$.
\item If
$\mathbb E_{\mathrm g}^{\mathbb S_b}([\bm U_b])
 =\mathbb E_{\mathrm g}^{\mathbb S_b}([\bm U'_b])$
in $R_g(\mathbb S_b)$ for every $b\in B$, then
$\mathbb E_{\mathrm g}^{\mathbb S}(\int[A])
 =\mathbb E_{\mathrm g}^{\mathbb S}(\int[A'])$ in $R_g(\mathbb S)$.
\end{enumerate}
\end{thm}

\begin{proof}
\emph{The bounded case.} Repeat the proof of Theorem~\ref{fubini} with the
presentations in place of the representatives. The finite package now
contains formulas for $\bm U_b$ and $\bm U'_b$ with their labels in
$\RV_\infty^{k}$ and $\RV_\infty^{l}$, finite-fibre bounds, the complete
graphs of both lift bijections, and the finite witness of the proof of
Theorem~\ref{fubini} for the equality of the bounded values of the
underlying RV sets: monomial charts, Euler certificates and a stabilized
bijection as in Lemma~\ref{test}. Validity is again first order, so compactness gives a finite
definable partition of $B$ with one package on each piece. The preparation
of the VF base retains all RV coordinates, including the vanishing labels.
The total RV object $\bm W$ has the structure map
$q(\tau,u)=(s,p_\tau(u))$, now with values in
$(\RV^\times)^d\times\RV_\infty^k$; its finiteness and the bijection
$\mathbb L\bm W\to Y$ are proved as before. The relative Gamma calculation
\eqref{relative:Gamma:push}, the relative torsor identity
\eqref{fubini:relative:torsor} and the stabilized witnesses concern only the
bounded values of RV families, which do not depend on the labels. The
cancellation of the proof of Theorem~\ref{fubini} therefore gives equality
of the bounded values of the two total objects. By
Lemma~\ref{lem:mixed:lift} these are the bounded values of the integrals of
the prepared total spaces. Finite additivity over the base partition, with
the value $0$ on empty pieces, proves (1).

\emph{The geometric case.} Fix $b$ temporarily. By
Lemma~\ref{lem:thickening}(2), $\theta\bm U_b\in\RV[k_b]$ and
$\theta\bm U'_b\in\RV[l_b]$ have the geometric values of $\bm U_b$ and
$\bm U'_b$, which are equal. By the denominator-clearing step
\eqref{geometric:fubini:clearing} there is $N\geq k_b,l_b$ such that the
unnormalized values of
$\theta\bm U_b\times\bm{\mathbb A}_{\RES}^{\,N-k_b}$ and
$\theta\bm U'_b\times\bm{\mathbb A}_{\RES}^{\,N-l_b}$ are equal in
$R(\mathbb S_b)$, and Lemma~\ref{test} supplies a finite stabilized witness
of this equality. The finite package consists of formulas for the two
presentations with finite-fibre bounds, the complete graphs of both lift
bijections, the integer $N$, and this witness for the padded thickenings.
By Lemma~\ref{lem:thickening}(1) the thickenings are given by fixed
formulas, with a fibre bound depending only on that of the presentation and
on one uniform bound for the label fibres of $\bm{\mathbb A}_{\RES}$.
Validity of a fixed package is first order over the original ground, and
the packages cover every ambient $b$. Compactness gives a finite definable
partition $(B_j)_j$ of $B$ with one package on each piece.

Work on one piece $B_j$, and write $\bm U$ and $\bm U'$ for the uniform
presentations. The lift graphs give $\mathbb S$-definable bijections of
$f^{-1}(B_j)$ and $(f')^{-1}(B_j)$ with $\mathbb L_{B_j}\bm U$ and
$\mathbb L_{B_j}\bm U'$. By Lemma~\ref{lem:thickening}(3), applied once for
every label,
\[
 \mathbb E_{\mathrm g}^{\mathbb S}\!\left(\int[\mathbb L_{B_j}\bm U]\right)
 =\mathbb E_{\mathrm g}^{\mathbb S}\!\left(\int[\mathbb L_{B_j}\theta\bm U]\right),
\]
and similarly for $\bm U'$. The thickened presentations have degrees
$k_b,l_b$ and nonzero labels. Together with the identity lift graphs, the
integer $N$ and the stabilized witness, they form a fixed single-degree
package on $B_j$. The argument after the compactness step in the proof of
Proposition~\ref{prop:geometric:fubini} applies to it and gives
\[
 \mathbb E_{\mathrm g}^{\mathbb S}\!\left(\int[\mathbb L_{B_j}\theta\bm U]\right)
 =\mathbb E_{\mathrm g}^{\mathbb S}\!\left(\int[\mathbb L_{B_j}\theta\bm U']\right).
\]
Combining these equalities, the restrictions of $A$ and $A'$ over $B_j$ have
equal geometric values. Finite additivity over the pieces, with the value
$0$ on empty pieces, proves (2).
\end{proof}

\begin{rem}
The single-degree hypotheses of Theorem~\ref{fubini} and
Proposition~\ref{prop:geometric:fubini} are the special cases of
Theorem~\ref{thm:mixed} with all labels nonzero. In the geometric case the
thickening is needed because a vanishing label changes the degree in which
its stratum is normalized; the bounded value is insensitive to labels.
Over the original ground, Lemma~\ref{lem:mixed:lift} shows that the value
\eqref{mixed:geometric:value} is the geometric value of the integral of the
lift. Hence the hypothesis of Theorem~\ref{thm:mixed}(2) is intrinsic
whenever $\mathbb S_b$ is itself a ground in the sense of
Section~\ref{sec:setup}, so that Lemma~\ref{lem:mixed:lift} applies over
it; the formalization proves that lemma over the original ground only.
\end{rem}
